Of course, it is not at all surprising that if a function is in $W^{1,p}(\Omega)$, its difference quotients obey some bound in terms of its partial derivatives. The following result is more substantial; it says that if we start out knowing that $u$ is in $L^p(\Omega)$ and can obtain a bound on its difference quotients that is independent of $h$, then we can deduce that $u$ is in the space $W^{1,p}(\Omega)$.

\begin{Lemma}
\label{lem:uniform-difference-quotient-converse}
Let $u\in L^p(\Omega)$, $1<p\leq\infty$, and suppose there exists a constant $\widetilde C$ such that for any $\Omega^{\prime}\subset\subset\Omega$ and $h<\operatorname{dist}(\Omega^{\prime},\partial\Omega)$ we have $D_i^hu\in L^p(\Omega^{\prime})$ and
\begin{equation}\label{eq:uniform-difference-quotient-converse-hypothesis}
\|D_i^hu\|_{L^p(\Omega^{\prime})}\leq\widetilde C.
\end{equation}
Then $u_{x_i}$ (which is a priori well defined as a distribution) is in fact in $L^p(\Omega)$ and satisfies
\begin{equation}\label{eq:uniform-difference-quotient-recovered-derivative-bound}
\|u_{x_i}\|_{L^p(\Omega)}\leq\widetilde C.
\end{equation}

\end{Lemma}

\begingroup
\renewcommand{\endProof}{\par\nobreak\hfill$\blacksquare$\par\vspace{0.1cm}}
\begin{Proof}
Recall that at the end of Section~\ref{subsec:boundary-value-problems-adjoints} we showed the distributional derivative of a function could be obtained as the limit of difference quotients. In terms of the present problem we have
\begin{equation}\label{eq:difference-quotient-distributional-derivative-limit}
(u_{x_i},\phi)=\lim_{h\to0}(D_i^hu,\phi)
\end{equation}
for each $\phi\in\mathcal{D}(\Omega)$.

We now note that there exists a function $v\in L^p(\Omega)$ with $\|v\|_{L^p(\Omega)}\leq\widetilde C$ and a sequence $h_m\to0$ such that for every $\phi\in\mathcal{D}(\Omega)$
\begin{equation}\label{eq:difference-quotient-weak-limit-function}
\lim_{h_m\to0}\int_\Omega\phi D_i^{h_m}u\,\mathrm{d}\bm{x}=\int_\Omega\phi v\,\mathrm{d}\bm{x}.
\end{equation}
This follows from the weak compactness of the bounded set $D_i^hu$ in $L^p(\Omega^{\prime})$ for every $\Omega^{\prime}\subset\subset\Omega$ (cf. Theorem~\ref{thm:alaoglu-weak-star-subsequence}) and the fact that $\Omega$ can be covered with a countable collection of closed subsets $\Omega_j\subset\subset\Omega$ such that each $\Omega^{\prime}\subset\subset\Omega$ intersects at most a finite number of the $\Omega_j$.

Thus,
\begin{equation}\label{eq:difference-quotient-recovered-distributional-derivative}
(u_{x_i},\phi)=\int_\Omega\phi v\,\mathrm{d}\bm{x};
\end{equation}
i.e., the distributional derivative of $u$ is given (uniquely) by the function $v\in L^p(\Omega)$.
\end{Proof}
\endgroup

We also need to develop a few important tools using difference operators; namely, the analogues of the product rule and integration by parts in differential calculus.

\begin{Lemma}
\label{lem:difference-quotient-product-and-parts}
Suppose $\Omega^{\prime}\subset\subset\Omega$ and $h<\operatorname{dist}(\Omega^{\prime},\partial\Omega)$. Then for any $u\in L^p(\Omega)$, $1<p<\infty$, and any $v\in C_0(\Omega)$ with $\operatorname{supp}v\subseteq\Omega^{\prime}$ we have
\begin{equation}\label{eq:difference-quotient-shifted-product-rule}
D_i^h(uv)(\bm{x})=u(\bm{x})D_i^hv(\bm{x})+D_i^hu(\bm{x})v(\bm{x}+h\bm{e}_i)
\end{equation}
and
\begin{equation}\label{eq:difference-quotient-integration-by-parts}
\int_\Omega D_i^hu(\bm{x})v(\bm{x})\,\mathrm{d}\bm{x}=-\int_\Omega u(\bm{x})D_i^{-h}v(\bm{x})\,\mathrm{d}\bm{x}.
\end{equation}

\end{Lemma}

The proof is left to the reader.
